\documentclass[aspectratio=169,12pt]{beamer}

% ── Theme ───────────────────────────────────────────────────────────────────
\usetheme{Madrid}
\usecolortheme{seahorse}
\setbeamertemplate{navigation symbols}{}
\setbeamertemplate{footline}[frame number]

% ── Packages ────────────────────────────────────────────────────────────────
\usepackage[utf8]{inputenc}
\usepackage[T5]{fontenc}
\usepackage{lmodern}
\usepackage[vietnamese]{babel}
\usepackage{booktabs}
\usepackage{multirow}
\usepackage{graphicx}
\usepackage{amsmath}
\usepackage{xcolor}
\usepackage{array}
\usepackage{colortbl}

% ── Colors ──────────────────────────────────────────────────────────────────
\definecolor{myblue}{RGB}{0,84,166}
\definecolor{mygreen}{RGB}{0,128,64}
\definecolor{myorange}{RGB}{204,85,0}
\definecolor{mygray}{RGB}{90,90,90}
\definecolor{highlight}{RGB}{255,240,0}

% ── Helpers ─────────────────────────────────────────────────────────────────
\newcommand{\good}[1]{\textcolor{mygreen}{\textbf{#1}}}
\newcommand{\bad}[1]{\textcolor{red}{\textbf{#1}}}
\newcommand{\sidenote}[1]{\textcolor{mygray}{\small #1}}

% ── Title ────────────────────────────────────────────────────────────────────
\title[Phát hiện DDoS bằng ML]{%
  Phát hiện Tấn công DDoS bằng Học máy:\\
  Pipeline FAMS + Leave-One-Group-Out\\
  trên CIC-DDoS 2019}
\subtitle{Báo cáo Học máy Ứng dụng}
\author{Nhóm nghiên cứu}
\institute{Học viện Kỹ thuật}
\date{Tháng 3, 2026}

% ════════════════════════════════════════════════════════════════════════════
\begin{document}
% ════════════════════════════════════════════════════════════════════════════

% ── Slide 1: Title ──────────────────────────────────────────────────────────
\begin{frame}
  \titlepage
\end{frame}

% ── Slide 2: Outline ────────────────────────────────────────────────────────
\begin{frame}{Nội dung trình bày}
  \tableofcontents[hideallsubsections]
\end{frame}

% ════════════════════════════════════════════════════════════════════════════
\section{Giới thiệu}
% ════════════════════════════════════════════════════════════════════════════

\begin{frame}{Tấn công DDoS là gì?}
  \begin{columns}[T]
    \column{0.55\textwidth}
      \begin{block}{Định nghĩa}
        Tấn công DDoS (Distributed Denial of Service) điều phối
        \textbf{nhiều nguồn phân tán} (botnet) để làm gián đoạn
        \textbf{tính sẵn sàng} của dịch vụ.
      \end{block}
      \vspace{0.3em}
      \begin{alertblock}{Hậu quả}
        \begin{itemize}
          \item Thiệt hại tài chính trực tiếp
          \item Mất uy tín tổ chức
          \item Gián đoạn hạ tầng thiết yếu
        \end{itemize}
      \end{alertblock}
    \column{0.42\textwidth}
      \begin{block}{Ba nhóm tấn công (CIC-DDoS 2019)}
        \small
        \begin{tabular}{lp{3.5cm}}
          \textbf{G1} & Reflection/Application\\
                      & DNS, LDAP, MSSQL,\\
                      & NetBIOS, SNMP, Portmap\\[4pt]
          \textbf{G2} & Reflection/Volume\\
                      & NTP, TFTP\\[4pt]
          \textbf{G3} & Exploitation/Flooding\\
                      & Syn, UDP, UDPLag,\\
                      & WebDDoS\\
        \end{tabular}
      \end{block}
  \end{columns}
\end{frame}

% ────────────────────────────────────────────────────────────────────────────
\begin{frame}{Thách thức phát hiện DDoS bằng học máy}
  \begin{enumerate}
    \item \textbf{Đa dạng loại tấn công} --- volumetric, protocol, application-layer\\
    \sidenote{Mỗi nhóm có phân phối lưu lượng rất khác nhau}
    \vspace{0.4em}
    \item \textbf{Khả năng tổng quát hóa} --- mô hình cần phát hiện tấn công \emph{mới chưa thấy}\\
    \sidenote{Câu hỏi trọng tâm của báo cáo này}
    \vspace{0.4em}
    \item \textbf{Mất cân bằng lớp} --- lưu lượng bình thường $\gg$ lưu lượng tấn công\\
    \sidenote{Dùng F2-score thay vì Accuracy}
    \vspace{0.4em}
    \item \textbf{Lựa chọn đặc trưng} --- 70+ đặc trưng, nhiều đặc trưng dư thừa\\
    \sidenote{Giải quyết bằng khung FAMS (Phase 1)}
  \end{enumerate}
\end{frame}

% ────────────────────────────────────────────────────────────────────────────
\begin{frame}{Câu hỏi nghiên cứu}
  \begin{block}{}
    \begin{enumerate}
      \item Tập đặc trưng FAMS từ bộ Kaggle có \textbf{hiệu quả trên CIC-DDoS 2019}?
      \vspace{0.5em}
      \item Mô hình huấn luyện trên một nhóm có \textbf{tổng quát hóa tốt} sang nhóm chưa thấy?
      \vspace{0.5em}
      \item Trong KNN, RF, XGB --- \textbf{mô hình nào tổng quát hóa tốt nhất} theo LOGO?
    \end{enumerate}
  \end{block}
  \vspace{0.5em}
  \begin{center}
    \textit{$\Rightarrow$ Trả lời qua thiết kế \textbf{Leave-One-Group-Out (LOGO)} trên CIC-DDoS 2019}
  \end{center}
\end{frame}

% ════════════════════════════════════════════════════════════════════════════
\section{Pipeline hai giai đoạn}
% ════════════════════════════════════════════════════════════════════════════

\begin{frame}{Tổng quan pipeline}
  \begin{columns}[T]
    \column{0.03\textwidth}
    \column{0.47\textwidth}
      \textcolor{myblue}{\textbf{Phase 1 --- Lựa chọn đặc trưng}}
      \vspace{0.3em}
      \begin{center}
        \fcolorbox{myblue}{blue!10}{\parbox{0.85\linewidth}{\centering\small Kaggle DDoS\\(CIC 2016/17/18)}}\\[0.3em]
        $\Downarrow$\\[0.3em]
        \fcolorbox{myblue}{blue!20}{\parbox{0.85\linewidth}{\centering\small FAMS\\(5 phương pháp + bỏ phiếu)}}\\[0.3em]
        $\Downarrow$\\[0.3em]
        \fcolorbox{myblue}{blue!35}{\parbox{0.85\linewidth}{\centering\small\textbf{18 đặc trưng}\\(votes $\geq 3$)}}
      \end{center}
    \column{0.03\textwidth}
      \vspace{2.8em}
      \begin{center}
        \Large $\Rightarrow$
      \end{center}
    \column{0.44\textwidth}
      \textcolor{myorange}{\textbf{Phase 2 --- Huấn luyện LOGO}}
      \vspace{0.3em}
      \begin{center}
        \fcolorbox{myorange}{orange!10}{\parbox{0.85\linewidth}{\centering\small CIC-DDoS 2019\\(17 Parquet)}}\\[0.3em]
        $\Downarrow$\\[0.3em]
        \fcolorbox{myorange}{orange!20}{\parbox{0.85\linewidth}{\centering\small LOGO Training\\(KNN / RF / XGB)}}\\[0.3em]
        $\Downarrow$\\[0.3em]
        \fcolorbox{myorange}{orange!35}{\parbox{0.85\linewidth}{\centering\small\textbf{Đánh giá}\\F2, Gap, FNR}}
      \end{center}
  \end{columns}
\end{frame}

% ════════════════════════════════════════════════════════════════════════════
\section{Phase 1 --- Lựa chọn đặc trưng FAMS}
% ════════════════════════════════════════════════════════════════════════════

\begin{frame}{Khung FAMS: bỏ phiếu đa phương pháp}
  \begin{columns}[T]
    \column{0.55\textwidth}
      \textbf{5 phương pháp lựa chọn đặc trưng:}
      \vspace{0.3em}
      \begin{itemize}
        \item \textcolor{myblue}{\textbf{Filter:}} Variance Threshold, Mutual Information
        \item \textcolor{myorange}{\textbf{Wrapper:}} RFE (Recursive Feature Elimination)
        \item \textcolor{mygreen}{\textbf{Embedded:}} Lasso L1, RF Feature Importance
      \end{itemize}
      \vspace{0.6em}
      \textbf{Cơ chế bỏ phiếu:}
      \[
        \mathcal{F}_{opt} = \left\{ f_j \;\middle|\; \sum_{k=1}^{5} \mathbf{1}[f_j \in \text{top-}25_k] \geq \tau \right\}
      \]

    \column{0.42\textwidth}
      \begin{block}{Hai ngưỡng so sánh}
        \centering
        \begin{tabular}{lcc}
          \toprule
          Ngưỡng $\tau$ & \#Đặc trưng & F2 TB\\
          \midrule
          $\geq 3$ & \good{18} & \good{0,9956}\\
          $> 3$    & 6          & 0,9840\\
          \bottomrule
        \end{tabular}
      \end{block}
      \vspace{0.4em}
      \begin{alertblock}{Lựa chọn}
        \textbf{18 đặc trưng (votes $\geq 3$)} được dùng cho Phase 2\\
        \small(giảm 75\% số chiều so với 70+ đặc trưng gốc)
      \end{alertblock}
  \end{columns}
\end{frame}

% ────────────────────────────────────────────────────────────────────────────
\begin{frame}{Phase 1: Kết quả đánh giá trên Kaggle}
  \begin{table}
    \centering
    \small
    \caption{Hiệu năng phân loại theo ngưỡng bỏ phiếu FAMS}
    \begin{tabular}{llrrrrr}
      \toprule
      \textbf{Ngưỡng} & \textbf{Mô hình} & \textbf{Acc} & \textbf{Prec} & \textbf{Recall} & \textbf{F2} & \textbf{FNR}\\
      \midrule
      \multirow{3}{*}{$\geq3$ (18 ft.)}
        & KNN & 0,993 & 0,992 & 0,994 & 0,993 & 0,006 \\
        & RF  & 0,997 & 0,998 & 0,997 & 0,997 & 0,003 \\
        & XGB & 0,997 & 0,998 & 0,996 & 0,997 & 0,004 \\
      \midrule
      \multirow{3}{*}{$>3$ (6 ft.)}
        & KNN & 0,983 & 0,978 & 0,987 & 0,985 & 0,013 \\
        & RF  & 0,982 & 0,981 & 0,981 & 0,981 & 0,019 \\
        & XGB & 0,983 & 0,977 & 0,988 & 0,986 & 0,012 \\
      \bottomrule
    \end{tabular}
  \end{table}
  \vspace{0.3em}
  \begin{center}
    \small FNR của RF giảm từ \bad{0,0188} xuống \good{0,0034} khi tăng từ 6 lên 18 đặc trưng
  \end{center}
\end{frame}

% ────────────────────────────────────────────────────────────────────────────
\begin{frame}{18 đặc trưng FAMS được chọn}
  \begin{columns}[T]
    \column{0.48\textwidth}
      \small
      \begin{enumerate}
        \item Fwd Packet Length Mean
        \item Fwd Packet Length Max
        \item Fwd Packets Length Total
        \item Fwd Seg Size Min
        \item Flow IAT Min
        \item Fwd Packet Length Std
        \item Subflow Bwd Bytes
        \item Packet Length Max
        \item Init Bwd Win Bytes
      \end{enumerate}
    \column{0.48\textwidth}
      \small
      \begin{enumerate}
        \setcounter{enumi}{9}
        \item Packet Length Std
        \item Subflow Fwd Bytes
        \item Fwd Header Length
        \item Flow Packets/s
        \item Bwd Packet Length Max
        \item Bwd Packets Length Total
        \item Avg Fwd Segment Size
        \item Subflow Fwd Packets
        \item Bwd Packet Length Std
      \end{enumerate}
  \end{columns}
  \vspace{0.4em}
  \begin{center}
    \sidenote{Chủ yếu: kích thước gói tin, tốc độ luồng, phân phối forward/backward}
  \end{center}
\end{frame}

% ════════════════════════════════════════════════════════════════════════════
\section{Phase 2 --- Thiết kế Leave-One-Group-Out}
% ════════════════════════════════════════════════════════════════════════════

\begin{frame}{CIC-DDoS 2019: Dữ liệu và phân nhóm}
  \begin{columns}[T]
    \column{0.55\textwidth}
      \begin{block}{Bộ dữ liệu}
        \begin{itemize}
          \item 17 file Parquet, 388.505 dòng (sau tiền xử lý)
          \item Đặc trưng: chỉ dùng \textbf{18 đặc trưng FAMS}
          \item Chuẩn hóa: \textbf{RobustScaler} (fit trên train)
        \end{itemize}
      \end{block}
      \vspace{0.4em}
      \begin{block}{Phân nhóm tấn công}
        \small
        \begin{tabular}{lll}
          \textbf{G1} & Refl/App & 22.839 mẫu\\
          \textbf{G2} & Refl/Vol & 200.705 mẫu\\
          \textbf{G3} & Exploit  & 77.471 mẫu\\
          \textbf{Benign} &      & 87.490 mẫu\\
        \end{tabular}
      \end{block}
    \column{0.42\textwidth}
      \begin{alertblock}{Thiết kế LOGO}
        \small
        \begin{itemize}
          \item \textbf{Exp 1:} Train G1 $\to$ Test G2+G3\\
                Train: 84.082 / Test: 304.423
          \vspace{0.2em}
          \item \textbf{Exp 2:} Train G2 $\to$ Test G1+G3\\
                Train: 261.948 / Test: 126.557
          \vspace{0.2em}
          \item \textbf{Exp 3:} Train G3 $\to$ Test G1+G2\\
                Train: 138.714 / Test: 249.791
        \end{itemize}
      \end{alertblock}
      \vspace{0.3em}
      \sidenote{Benign 70\% cố định cho train, 30\% cho test --- dùng chung ba Exp}
  \end{columns}
\end{frame}

% ────────────────────────────────────────────────────────────────────────────
\begin{frame}{Quy trình huấn luyện và tối ưu siêu tham số}
  \begin{enumerate}
    \item \textbf{Optuna TPE} --- 30 trials, mục tiêu tối đa \textbf{F2 (5-fold StratifiedKFold)}\\
    \sidenote{Không gian tìm kiếm: $k$, metric (KNN); n\_estimators, depth, features (RF/XGB)}
    \vspace{0.4em}
    \item \textbf{Xử lý mất cân bằng lớp}
    \begin{itemize}
      \item RF: \texttt{class\_weight='balanced'}
      \item XGB: \texttt{scale\_pos\_weight}
      \item KNN: không hỗ trợ tường minh
    \end{itemize}
    \vspace{0.4em}
    \item \textbf{Huấn luyện lại} trên toàn bộ train sau khi tìm được best params
    \vspace{0.4em}
    \item \textbf{Đánh giá} trên tập test (nhóm chưa thấy trong huấn luyện)
  \end{enumerate}
  \vspace{0.5em}
  \begin{block}{Chỉ số đánh giá chính}
    \centering
    \textbf{F2-score}, F1, Recall, FNR, AUC, \textbf{Gap} = CV F2 $-$ Test F2
  \end{block}
\end{frame}

% ════════════════════════════════════════════════════════════════════════════
\section{Kết quả thực nghiệm}
% ════════════════════════════════════════════════════════════════════════════

\begin{frame}{Kết quả tổng hợp: Leave-One-Group-Out}
  \begin{table}
    \centering
    \small
    \begin{tabular}{llrrrrr}
      \toprule
      \textbf{Thí nghiệm} & \textbf{Mô hình} & \textbf{CV F2} & \textbf{Test F2} & \textbf{Recall} & \textbf{FNR} & \textbf{Gap}\\
      \midrule
      \multirow{3}{*}{Exp 1 (G1$\to$G2+G3)}
        & KNN  & 0,9952 & \bad{0,603}  & 0,548 & 0,452 & +0,39 \\
        & RF   & 0,9972 & \good{0,833} & 0,800 & 0,200 & +0,16 \\
        & XGB  & 0,9976 & 0,776        & 0,735 & 0,265 & +0,22 \\
      \midrule
      \multirow{3}{*}{Exp 2 (G2$\to$G1+G3)}
        & KNN  & 0,9991 & \bad{0,622}  & 0,568 & 0,432 & +0,38 \\
        & RF   & 0,9995 & \bad{0,620}  & 0,566 & 0,434 & +0,38 \\
        & XGB  & 0,9993 & \bad{0,622}  & 0,568 & 0,432 & +0,38 \\
      \midrule
      \multirow{3}{*}{Exp 3 (G3$\to$G1+G2)}
        & KNN  & 0,9937 & 0,917        & 0,899 & 0,101 & +0,08 \\
        & \textbf{RF}  & \textbf{0,9981} & \good{\textbf{0,985}} & \textbf{0,981} & \good{\textbf{0,019}} & \good{\textbf{+0,01}} \\
        & XGB  & 0,9978 & 0,902        & 0,881 & 0,119 & +0,10 \\
      \bottomrule
    \end{tabular}
  \end{table}
  \vspace{-0.2em}
  \begin{center}
    \small \textbf{RF Exp 3}: Test F2 = 0,985; Gap = +0,013; FNR = \good{1,9\%}
  \end{center}
\end{frame}

% ────────────────────────────────────────────────────────────────────────────
\begin{frame}{Phân tích: Exp 1 --- Huấn luyện trên G1}
  \begin{columns}[T]
    \column{0.5\textwidth}
      \begin{block}{Bối cảnh}
        \begin{itemize}
          \item G1 nhỏ nhất: 22.839 mẫu tấn công
          \item Đặc trưng đồng nhất: phản xạ giao thức ứng dụng (DNS/LDAP/MSSQL)
        \end{itemize}
      \end{block}
      \vspace{0.3em}
      \begin{alertblock}{Kết quả}
        \begin{itemize}
          \item KNN: Test F2 = \bad{0,60}, Gap +0,39
          \item RF:  Test F2 = \good{0,83}, Gap +0,16
          \item XGB: Test F2 = 0,78, Gap +0,22
        \end{itemize}
      \end{alertblock}
    \column{0.47\textwidth}
      \begin{block}{Giải thích}
        Mô hình học ``nhận dạng giao thức cụ thể'' thay vì pattern DDoS tổng quát.
        RF tổng quát hóa tốt hơn nhờ \textbf{bagging} đa dạng hóa biên quyết định.
      \end{block}
  \end{columns}
\end{frame}

% ────────────────────────────────────────────────────────────────────────────
\begin{frame}{Phân tích: Exp 2 --- Huấn luyện trên G2}
  \begin{columns}[T]
    \column{0.5\textwidth}
      \begin{block}{Bối cảnh}
        \begin{itemize}
          \item G2 lớn nhất: 200.705 mẫu
          \item Nhưng \textbf{rất đồng nhất}: chỉ NTP (UDP/123) và TFTP (UDP/69)
        \end{itemize}
      \end{block}
      \vspace{0.3em}
      \begin{alertblock}{Kết quả đáng chú ý}
        \begin{itemize}
          \item Cả 3 mô hình: Test F2 $\approx$ \bad{0,62}
          \item Gap đồng loạt +0,38
          \item \textbf{Lợi thế bagging của RF biến mất!}
        \end{itemize}
      \end{alertblock}
    \column{0.47\textwidth}
      \begin{block}{Bài học quan trọng}
        \textit{Kích thước tập train lớn \textbf{không đảm bảo} tổng quát hóa tốt nếu tập đó \textbf{thiếu đa dạng} về loại lưu lượng.}
        \vspace{0.4em}

        Khi toàn bộ 200k mẫu phân bố trong vùng đặc trưng hẹp (NTP/TFTP), mọi cây trong RF đều học cùng biên quyết định hẹp.
      \end{block}
  \end{columns}
\end{frame}

% ────────────────────────────────────────────────────────────────────────────
\begin{frame}{Phân tích: Exp 3 --- Huấn luyện trên G3 (\textit{Best case})}
  \begin{columns}[T]
    \column{0.5\textwidth}
      \begin{block}{Bối cảnh}
        \begin{itemize}
          \item G3 đa dạng: Syn, UDP, UDPLag, WebDDoS
          \item Bao gồm cả tấn công \textbf{thể tích} lẫn \textbf{tầng ứng dụng}
          \item Tương đồng với dữ liệu Kaggle (Phase 1)
        \end{itemize}
      \end{block}
      \vspace{0.3em}
      \begin{block}{Kết quả xuất sắc}
        \begin{itemize}
          \item RF: Test F2 = \good{\textbf{0,985}}, Gap = \good{+0,013}
          \item FNR = \good{1,89\%} --- phát hiện hơn 98\% tấn công từ G1+G2
        \end{itemize}
      \end{block}
    \column{0.47\textwidth}
      \begin{alertblock}{Lý giải}
        Đặc trưng học từ G3 \textbf{đa dạng về kích thước gói, tốc độ luồng, tỷ lệ forward/backward} → tổng quát hóa tốt sang G1+G2.

        \vspace{0.4em}
        FAMS Phase 1 cũng dùng Kaggle chứa UDP/Syn Flood → tập 18 đặc trưng phù hợp với G3.
      \end{alertblock}
  \end{columns}
\end{frame}

% ────────────────────────────────────────────────────────────────────────────
\begin{frame}{Hiệu quả tính toán}
  \begin{table}
    \centering
    \small
    \caption{Thời gian huấn luyện và suy luận (KNN dùng \texttt{algorithm=ball\_tree})}
    \begin{tabular}{llrrrr}
      \toprule
      \textbf{Thí nghiệm} & \textbf{Mô hình} & \textbf{Train size} & \textbf{Train (s)} & \textbf{Infer (ms/mẫu)} & \textbf{Test F2}\\
      \midrule
      \multirow{3}{*}{Exp 1 (84k)}
        & KNN  & 84.082  & 0,27  & 0,3148 & 0,603 \\
        & RF   & 84.082  & 5,85  & 0,0161 & 0,833 \\
        & XGB  & 84.082  & 1,97  & \good{0,0034} & 0,776 \\
      \midrule
      \multirow{3}{*}{Exp 2 (262k)}
        & KNN  & 261.948 & 1,51  & 0,1760 & 0,622 \\
        & RF   & 261.948 & 1,18  & \good{0,0027} & 0,620 \\
        & XGB  & 261.948 & 13,11 & 0,0082 & 0,622 \\
      \midrule
      \multirow{3}{*}{Exp 3 (139k)}
        & KNN  & 138.714 & 0,08  & 0,6512 & 0,917 \\
        & RF   & 138.714 & 4,94  & 0,0106 & \good{0,985} \\
        & XGB  & 138.714 & 4,28  & \good{0,0011} & 0,902 \\
      \bottomrule
    \end{tabular}
  \end{table}
  \vspace{0.2em}
  \sidenote{RF và XGB: thời gian suy luận $<$ 0,02 ms/mẫu --- phù hợp phát hiện thời gian thực}
\end{frame}

% ════════════════════════════════════════════════════════════════════════════
\section{Thảo luận}
% ════════════════════════════════════════════════════════════════════════════

\begin{frame}{So sánh ba mô hình}
  \begin{columns}[T]
    \column{0.32\textwidth}
      \begin{block}{\centering Random Forest}
        \centering
        \good{Tốt nhất tổng quát}\\
        \vspace{0.3em}
        \small
        Test F2 cao nhất ở Exp 1 (0,83) và Exp 3 (0,985).\\
        Gap nhỏ nhất ở Exp 3 (+0,013).\\
        Bagging giảm phương sai hiệu quả.
      \end{block}
    \column{0.32\textwidth}
      \begin{block}{\centering XGBoost}
        \centering
        \textcolor{myorange}{AUC cao}\\
        \vspace{0.3em}
        \small
        AUC = 0,99 ở Exp 1, Exp 3.\\
        Nhưng Gap lớn hơn RF ở Exp 3 (+0,096 vs +0,013).\\
        Học pattern phức tạp, dễ overfit.
      \end{block}
    \column{0.32\textwidth}
      \begin{block}{\centering KNN}
        \centering
        \bad{Phụ thuộc phân phối}\\
        \vspace{0.3em}
        \small
        Test F2 thấp nhất ở Exp 1, 2 ($\approx$ 0,61).\\
        Cải thiện tốt khi train/test tương đồng (Exp 3: 0,92).\\
        Suy luận chậm hơn RF/XGB.
      \end{block}
  \end{columns}
\end{frame}

% ────────────────────────────────────────────────────────────────────────────
\begin{frame}{Trả lời câu hỏi nghiên cứu}
  \begin{block}{Q1: FAMS Kaggle có hiệu quả trên CIC-DDoS 2019?}
    \textit{Có, nhưng không đồng đều.} Test F2 = 0,985 (RF, Exp 3) khi nhóm train tương đồng với Kaggle. Giảm rõ rệt (0,62--0,83) khi nhóm train là Reflection.
  \end{block}
  \vspace{0.3em}
  \begin{block}{Q2: Tổng quát hóa sang nhóm chưa thấy?}
    \textit{Phụ thuộc tính đa dạng của nhóm train.} G3 (đa dạng) → Gap $\approx$ 0,01. G1, G2 (đơn điệu giao thức) → Gap = 0,16--0,39.
  \end{block}
  \vspace{0.3em}
  \begin{block}{Q3: Mô hình nào tốt nhất theo LOGO?}
    \textbf{Random Forest} --- nhất quán Test F2 và Gap tốt nhất qua cả ba thí nghiệm.
  \end{block}
\end{frame}

% ────────────────────────────────────────────────────────────────────────────
\begin{frame}{Phát hiện quan trọng: Đa dạng $>$ Kích thước}
  \begin{center}
    \Large
    \textbf{Tính đa dạng của dữ liệu huấn luyện}\\
    --- không phải kích thước ---\\
    \textbf{là yếu tố quyết định khả năng tổng quát hóa}
  \end{center}
  \vspace{0.5em}
  \begin{columns}[T]
    \column{0.47\textwidth}
      \begin{alertblock}{Bằng chứng}
        G2 (200k mẫu, nhưng chỉ NTP+TFTP):\\
        Test F2 = \bad{0,62}\\[0.3em]
        G3 (78k mẫu, nhưng Syn+UDP+WebDDoS):\\
        Test F2 = \good{0,985}
      \end{alertblock}
    \column{0.47\textwidth}
      \begin{block}{Ý nghĩa thực tiễn}
        Chiến lược thu thập dữ liệu huấn luyện nên ưu tiên \textbf{đa dạng loại tấn công} hơn \textbf{số lượng mẫu thuần túy}.
      \end{block}
  \end{columns}
\end{frame}

% ════════════════════════════════════════════════════════════════════════════
\section{Kết luận}
% ════════════════════════════════════════════════════════════════════════════

\begin{frame}{Kết luận}
  \begin{block}{Đóng góp chính}
    \begin{enumerate}
      \item \textbf{Lựa chọn đặc trưng FAMS đa phương pháp:} 18 đặc trưng (votes $\geq 3$) đạt F2 TB = 0,9956 trên Kaggle; giảm 75\% số chiều.

      \vspace{0.4em}
      \item \textbf{Đánh giá LOGO trên CIC-DDoS 2019:}
        \begin{itemize}
          \item RF Exp 3: Test F2 = 0,985, FNR = 1,89\%, Gap $\approx$ 0 --- không overfit nhóm
          \item Exp 1\&2: Gap lớn (0,16--0,39) do đặc trưng giao thức đặc thù không chuyển giao tốt
          \item Random Forest nhất quán tốt nhất trong ba mô hình
        \end{itemize}
    \end{enumerate}
  \end{block}
  \vspace{0.3em}
  \begin{alertblock}{Phát hiện cốt lõi}
    \textbf{Đa dạng dữ liệu huấn luyện $>$ Kích thước} --- quyết định khả năng tổng quát hóa
  \end{alertblock}
\end{frame}

% ────────────────────────────────────────────────────────────────────────────
\begin{frame}{Hạn chế và hướng phát triển}
  \begin{columns}[T]
    \column{0.48\textwidth}
      \begin{alertblock}{Hạn chế}
        \begin{itemize}
          \item Dữ liệu phòng thí nghiệm (chưa có TLS, IPv6)
          \item Bất đối xứng tổng quát hóa: tốt từ G3 sang G1+G2, kém chiều ngược lại
          \item Mô hình tĩnh, chưa xử lý \textit{concept drift}
          \item WebDDoS chỉ 51 mẫu trong G3
        \end{itemize}
      \end{alertblock}
    \column{0.48\textwidth}
      \begin{block}{Hướng phát triển}
        \begin{itemize}
          \item Giải quyết bất đối xứng tổng quát hóa (transfer learning)
          \item Huấn luyện đa nhóm với tỷ lệ mẫu tối ưu
          \item Đánh giá trên dữ liệu thực (TLS, cloud)
          \item Online Learning / periodic retraining
          \item SHAP để giải thích feature importance
        \end{itemize}
      \end{block}
  \end{columns}
\end{frame}

% ────────────────────────────────────────────────────────────────────────────
\begin{frame}
  \begin{center}
    \vspace{1em}
    {\Large\textbf{Cảm ơn!}}\\[1em]
    \normalsize
    Code: \texttt{phase1\_feature\_selection.ipynb}\\
    \texttt{phase2\_model\_training.ipynb}\\[1em]
    {\large Câu hỏi và thảo luận?}
  \end{center}
\end{frame}

% ════════════════════════════════════════════════════════════════════════════
% Backup slides
% ════════════════════════════════════════════════════════════════════════════
\appendix

\begin{frame}{Siêu tham số tốt nhất (Optuna TPE, 30 trials)}
  \begin{table}
    \centering
    \small
    \begin{tabular}{llp{7.5cm}}
      \toprule
      \textbf{Exp} & \textbf{Mô hình} & \textbf{Siêu tham số tốt nhất}\\
      \midrule
      \multirow{3}{*}{Exp 1}
        & KNN  & \texttt{n\_neighbors=3, weights=distance} \\
        & RF   & \texttt{n\_estimators=100, max\_depth=19, min\_samples\_split=2, max\_features=log2} \\
        & XGB  & \texttt{n\_estimators=500, max\_depth=7, lr=0.035, subsample=0.785} \\
      \midrule
      \multirow{3}{*}{Exp 2}
        & KNN  & \texttt{n\_neighbors=5, weights=distance} \\
        & RF   & \texttt{n\_estimators=150, max\_depth=18, min\_samples\_split=2, max\_features=log2} \\
        & XGB  & \texttt{n\_estimators=450, max\_depth=7, lr=0.094, subsample=0.999} \\
      \midrule
      \multirow{3}{*}{Exp 3}
        & KNN  & \texttt{n\_neighbors=3, weights=distance} \\
        & RF   & \texttt{n\_estimators=500, max\_depth=15, min\_samples\_split=5, max\_features=log2} \\
        & XGB  & \texttt{n\_estimators=200, max\_depth=6, lr=0.133, subsample=0.918} \\
      \bottomrule
    \end{tabular}
  \end{table}
  \sidenote{KNN cố định \texttt{algorithm=ball\_tree} cho mọi thí nghiệm}
\end{frame}

\begin{frame}{F$\beta$-score: tại sao chọn F2?}
  \[
    F_\beta = (1+\beta^2)\, \frac{\text{Precision} \times \text{Recall}}{\beta^2 \cdot \text{Precision} + \text{Recall}}
  \]
  \begin{itemize}
    \item $\beta = 1$ (F1): cân bằng Precision và Recall
    \item $\beta = 2$ (F2): \textbf{ưu tiên Recall gấp đôi Precision}
  \end{itemize}
  \vspace{0.5em}
  \begin{block}{Lý do chọn F2 trong phát hiện DDoS}
    Bỏ sót tấn công (False Negative, FNR cao) → hậu quả nghiêm trọng hơn\\
    báo động nhầm (False Positive).\\
    F2 làm mục tiêu Optuna → khuyến khích Recall cao.
  \end{block}
\end{frame}

\begin{frame}{RobustScaler: tại sao không dùng StandardScaler?}
  \[
    x_{\text{RobustScaler}} = \frac{x - Q_2(x)}{Q_3(x) - Q_1(x)}
  \]
  \begin{columns}[T]
    \column{0.48\textwidth}
      \begin{block}{StandardScaler}
        Dùng mean và std → \textbf{nhạy với ngoại lai}\\
        Lưu lượng burst traffic có thể làm lệch mean đáng kể
      \end{block}
    \column{0.48\textwidth}
      \begin{block}{RobustScaler}
        Dùng trung vị ($Q_2$) và IQR ($Q_3-Q_1$) → \textbf{bền với ngoại lai}\\
        Phù hợp dữ liệu lưu lượng mạng
      \end{block}
  \end{columns}
  \vspace{0.4em}
  \begin{alertblock}{Quan trọng}
    Scaler được \textit{fit} chỉ trên tập train → transform cả train và test\\
    $\Rightarrow$ tránh data leakage
  \end{alertblock}
\end{frame}

\end{document}
